% TOP_PROBLEM: {"id": 30007502, "title": "Filtered chromatic splitting at generic primes", "category_id": 7, "category": {"id": 7, "name": "topology", "display_name": "Topology", "description": "Properties preserved under continuous deformations.", "slug": "topology", "order_index": 7, "created_at": "2026-07-31T15:26:25.671Z"}, "status": "open"}
% FIRST_POSTED: 2026-10-06
% ATTEMPT_STATUS: solved
% OPENAI_CLAIM: {"schema_version":1,"source_repo":"https://github.com/openai/math","source_commit":"adc7f1241b42e322a6451854ab7e4b4c146bf78a","release_date":"2026-10-06","match":"full","resolution":"proved","independently_reviewed":false,"second_pass":"agreed","families":[{"family":"318","title":"Chromatic splitting: filtrations and counterexamples","subject":"Topology","manuscripts":[{"title":"Filtered chromatic splitting at generic primes","pdf_path":"preprints/Filtered-chromatic-splitting-at-generic-primes-September-25-2026/paper.pdf","readme_path":"preprints/Filtered-chromatic-splitting-at-generic-primes-September-25-2026/README.md","date":"2026-09-25","theorem_ref":"thm:main (sections/01-introduction.tex:33)","statement_sources":[{"path":"preprints/Filtered-chromatic-splitting-at-generic-primes-September-25-2026/build/sections/01-introduction.tex","line":33,"label":"thm:main"}]}],"lean":null,"reasoning_trace":null}]}
\documentclass[11pt]{article}
\usepackage[margin=25mm]{geometry}
\usepackage{fontspec}
\setmainfont{Latin Modern Roman}
\usepackage{amsmath,amssymb,mathtools}
\usepackage{unicode-math}
\setmathfont{Latin Modern Math}
\usepackage{enumitem}
\usepackage{xurl,hyperref,bookmark}
\hypersetup{colorlinks=true,urlcolor=blue}
\setcounter{secnumdepth}{0}
\setlength{\parindent}{0pt}
\setlength{\parskip}{5pt}
\setlength{\emergencystretch}{3em}
\begin{document}
\section*{Filtered chromatic splitting at generic primes}

\subsection{OpenAI's claim}
For n at least one and p greater than n+1, lower-height localization of the K(n)-local sphere is claimed to have an ordered \(2^n\)-stage filtration with the classical chromatic-splitting cofibers and canonical initial unit. The cited manuscript argues for that result; its relation to this problem is specified below.

\subsection{Mathematical statement of the claimed result}
\paragraph{Mathematical claim 1.} For $n\ge1$ and prime $p>n+1$, put $S=S_p^\wedge$, $D=L_{K(n)}S$ and $X_{n,p}=L_{n-1}D$, where $L_j=L_{E(j)}$. OpenAI claims a filtration in $E(n-1)$-local $S$-modules with $2^n$ cofibers indexed by subsets $I\subseteq\{1,\ldots,n\}$:
\[
C_\varnothing=L_{n-1}S,\qquad C_I=\Sigma^{-d(I)}L_{n-\max I}S\ (I\ne\varnothing),\qquad d(I)=\sum_{i\in I}(2i-1).
\]
The empty subset is first, the remaining maxima are nondecreasing, and the first-layer map is the canonical unit. It claims a single family $y_i\in\pi_{1-2i}D$ whose ordered products $y_I$ give, for every $1\le t<n$, equivalences
\[
L_{K(t)}\left(\bigvee_{I\subseteq\{1,\ldots,n-t\}}\Sigma^{-d(I)}S\xrightarrow{(y_I)}D\right).
\]
A filtered extension is asserted, not a split wedge decomposition before localization.

\subsection{Outline of the argument}
Manuscript 1 (family 318, 2026-09-25): Finite marking and transfer calculations first establish coefficient finiteness. A separately constructed integral family of sphere-cochain classes has normalized exterior products acting as bases under lower-height tests. Removing blocks in decreasing chromatic height yields the ordered cofibers; independent rational dimension counts supply the final block while retaining the extensions between stages.

\subsection{Scope relative to this problem}
Manuscript 1 (family 318, 2026-09-25): The problem selects the mathematical statement summarized here, including its stated hypotheses and parameter range. The cited theorem gives a claimed proof of that assertion in full.

\subsection{Formal verification}
Family 318 has no Lean scope page in this snapshot; no formalization of these manuscript results is listed.

This site has not executed the Lean files or checked their dependencies. These descriptions report the scope of the linked declarations, rather than a new verification of the proofs.

\subsection{Status caveat}
OpenAI attributes these results to an internal model. Its release notice cautions that results lacking full formalization can contain mistakes. This site has not independently assessed the proofs; a solved label records the scope of the reported claim.

\subsection{Sources}
\begin{itemize}
\item OpenAI, \emph{Filtered chromatic splitting at generic primes}, 2026-09-25; citation key \texttt{OAI:Filtered-chromatic-splitting-at-generic-primes-September-25-2026}. \href{https://github.com/openai/math/blob/adc7f1241b42e322a6451854ab7e4b4c146bf78a/preprints/Filtered-chromatic-splitting-at-generic-primes-September-25-2026/paper.pdf}{PDF (GitHub)}; \href{https://github.com/openai/math/raw/adc7f1241b42e322a6451854ab7e4b4c146bf78a/preprints/Filtered-chromatic-splitting-at-generic-primes-September-25-2026/paper.pdf}{PDF (direct)}; \href{https://github.com/openai/math/blob/adc7f1241b42e322a6451854ab7e4b4c146bf78a/preprints/Filtered-chromatic-splitting-at-generic-primes-September-25-2026/README.md}{Manuscript page and citation}; \href{https://github.com/openai/math/blob/main/preprints/Filtered-chromatic-splitting-at-generic-primes-September-25-2026/paper.pdf}{Latest version}; Statement source: \href{https://github.com/openai/math/blob/adc7f1241b42e322a6451854ab7e4b4c146bf78a/preprints/Filtered-chromatic-splitting-at-generic-primes-September-25-2026/build/sections/01-introduction.tex\#L33}{thm:main}.
\item \href{https://github.com/openai/math/blob/adc7f1241b42e322a6451854ab7e4b4c146bf78a/CONTENTS.md}{OpenAI catalogue entry}.
\end{itemize}
\end{document}
